\documentclass[11pt]{article}

\usepackage[margin=0.8in]{geometry}
\usepackage{booktabs}
\usepackage{array}
\usepackage{caption}
\usepackage[T1]{fontenc}
\usepackage{lmodern}

\captionsetup{font=small,labelfont=bf}

\begin{document}

\section*{SERS Substrate-Agnostic Dataset Matrix}

The current substrate-agnostic evaluation uses canonical chemical labels. In
this analysis, BT is merged into benzenethiol. The table below shows the current
number of spectra for each chemical-substrate pair.

\begin{table}[h]
\centering
\caption{Current canonical three-chemical SERS matrix.}
\begin{tabular}{lrrrrrrr}
\toprule
Chemical & Ag & AgNP & Au & AuNP & PICO & pSERS & Current total \\
\midrule
\texttt{4np} & 0 & 25 & 0 & 0 & 25 & 25 & 75 \\
\texttt{benzenethiol} & 25 & 0 & 25 & 0 & 25 & 25 & 100 \\
\texttt{pyridine} & 0 & 25 & 0 & 25 & 25 & 25 & 100 \\
\bottomrule
\end{tabular}
\end{table}

The minimum recommended expansion is to fill every missing chemical-substrate
cell before heavily increasing spectra from already-covered cells.

\begin{table}[h]
\centering
\caption{Minimum target matrix for stronger substrate-agnostic evaluation.}
\begin{tabular}{lrrrrrrr}
\toprule
Chemical & Ag & AgNP & Au & AuNP & PICO & pSERS & Target total \\
\midrule
\texttt{4np} & 25+ & 25+ & 25+ & 25+ & 25+ & 25+ & 150+ \\
\texttt{benzenethiol} & 25+ & 25+ & 25+ & 25+ & 25+ & 25+ & 150+ \\
\texttt{pyridine} & 25+ & 25+ & 25+ & 25+ & 25+ & 25+ & 150+ \\
\bottomrule
\end{tabular}
\end{table}

\noindent
If time allows, the stronger target is two to three independent
preparations/maps per chemical-substrate pair, with approximately 25 spectra per
preparation. Independent preparations are more useful for substrate-agnostic
validation than many correlated spectra from a single map.

\end{document}
